% BUILD11. Descritivo.
\begin{table}[htbp]\centering
\caption{Extreme-case imputations for non-entrants in the decomposition of the fall in age at entry, birth cohorts 2010--2014}
\label{tab:E5}
\footnotesize
\setlength{\tabcolsep}{4pt}
\begin{tabular}{lrrrr}
\toprule
Scenario & $\Delta$ appl.\ (mo.) & $\Delta$ entry (mo.) & Share appl. & Share wait \\
\midrule
Enrolled only (paper) & -8.68 & -9.37 & 93\% & 7\% \\
Combined, aged-out at 48 (max entry) & -14.43 & -12.56 & 115\% & -15\% \\
Combined, aged-out at appl.\ age (min) & -14.43 & -12.85 & 112\% & -12\% \\
\bottomrule
\end{tabular}
\\[3pt]{\footnotesize \textit{Note:} Age at entry is undefined for a child who never enters, so the full-cohort decomposition is not an estimand these data identify; this table reports the split under the two extreme imputations of that missing entry age rather than a bound. Children who aged out of eligibility unreferred (27.4 percent of the 2010 cohort, 5.2 percent of 2014) are assigned, as a worst case, an entry age equal to their age at application (no wait) or, as the other extreme, equal to the 48-month limit. Row 1 is the paper's enrolled-only figure. Age at application is observed for all, so its change ($\Delta$ appl.) is fixed in the combined population; only the imputed entry-age change moves the split. A share above 100 percent means the implied wait component is \emph{negative}---waiting rises once the never-entering children are counted.}
\end{table}
